%-------------------------
% Alexander Tong | Space Hardware Resume
% Adapted from Jake Gutierrez's resume template
% Compile with pdfLaTeX in Overleaf
%-------------------------

\documentclass[letterpaper,11pt]{article}

\usepackage[T1]{fontenc}
\usepackage[utf8]{inputenc}
\usepackage{lmodern}
\usepackage[margin=0.5in]{geometry}
\usepackage{titlesec}
\usepackage{enumitem}
\usepackage[hidelinks]{hyperref}
\usepackage[english]{babel}
\input{glyphtounicode}

\pdfgentounicode=1
\pagestyle{empty}
\urlstyle{same}
\raggedright
\raggedbottom
\setlength{\parindent}{0pt}
\setlength{\tabcolsep}{0pt}

% Section formatting
\titleformat{\section}
  {\large\scshape\raggedright}
  {}{0em}{}[\titlerule]
\titlespacing*{\section}{0pt}{8pt}{5pt}

% Single-column headings preserve a straightforward reading order.
\newcommand{\resumeSubheading}[4]{
  \item
  {\small\textbf{#1}\hfill #2}\\
  {\small\textit{#3}\hfill\textit{#4}}
}

\newcommand{\resumeItem}[1]{
  \item \small{#1}
}

\newcommand{\resumeSubHeadingListStart}{
  \begin{itemize}[
    leftmargin=0pt,
    label={},
    itemsep=4pt,
    topsep=0pt,
    parsep=0pt
  ]
}

\newcommand{\resumeSubHeadingListEnd}{
  \end{itemize}
}

\newcommand{\resumeItemListStart}{
  \begin{itemize}[
    leftmargin=0.16in,
    label=\textbullet,
    itemsep=1pt,
    topsep=2pt,
    parsep=0pt,
    partopsep=0pt
  ]
}

\newcommand{\resumeItemListEnd}{
  \end{itemize}
}

\begin{document}

%----------HEADER----------
\begin{center}
  {\Huge\bfseries\scshape Alexander Tong}\\[4pt]
  \small
  860-706-3624 \enspace\textbar\enspace
  \href{mailto:alto335@bu.edu}{alto335@bu.edu}
  \enspace\textbar\enspace
  \href{https://alexandertong.space}{alexandertong.space}
  \enspace\textbar\enspace
  \href{https://www.linkedin.com/in/alto335/}{linkedin.com/in/alto335}
\end{center}
\vspace{-9pt}

%----------EDUCATION----------
\section{Education}
{\small
\textbf{Boston University}\hfill Expected May 2028\\
B.S. Mechanical Engineering, Aerospace Concentration
\hfill GPA: 3.71/4.00
}

\section{Experience}
\resumeSubHeadingListStart
\resumeSubheading
  {Boston University Plummer Laboratory}
  {June 2026 - Present}
  {Mechanical Engineering Intern}
  {Boston, MA}
\resumeItemListStart
  \resumeItem{Designed a dynamic-test fixture using SolidWorks FEA, predicting 0.44 mm deflection under a 3 kg load at 20.1 Hz}
  \resumeItem{Calculated bolt preload and joint failure margins: 1.49 FOS for peak fastener loading and minimum 2.46 FOS for thread shear, plate bearing, and pull-through}
  \resumeItem{Created engineering drawings, G-code, GD\&T, and worst-case tolerance stacks; CNC-machined a 6061-T6 adapter on a Haas VF-2 with 12 drilled holes held to $\pm$0.005 in}
  \resumeItem{Designed the test matrix, procedure, and three of four specimen concepts for 48 Instron tests; processed 0--181 mN force data}
\resumeItemListEnd

\resumeSubheading
  {Culturon}
  {February 2026 - June 2026}
  {Mechanical Engineering Intern}
  {Sydney, Australia}
\resumeItemListStart
  \resumeItem{Planned and executed the company's first full-scale nanoparticle deposition mapping across chamber locations and magnetic configurations using SEM, optical microscopy, visual evidence, and ImageJ}
  \resumeItem{Formed a working hypothesis for contamination cause and pattern from sample-plate results, microscopy, and visual evidence}
  \resumeItem{Presented pump-port/chamber and electron-curtain prevention concepts based on the hypothesis and fluid/electromagnetic first principles in a technical design review}
  \resumeItem{Designed Onshape sheet-metal sample racks for reactor-port insertion and adhesive-free mounting}
\resumeItemListEnd

\resumeSubHeadingListEnd
\section{Engineering Projects}
\resumeSubHeadingListStart
\resumeSubheading
  {Argo Separation and Recovery Subsystem}
  {June 2025 - Present}
  {Co-Lead Engineer, Boston University Rocket Propulsion Group}
  {}
\resumeItemListStart
  \resumeItem{Designed redundant recovery for a 135 lb dry-mass rocket targeting 100,000 ft through four design reviews; final subsystem mass 18.7 lbm within a 20 lbm budget, final BOM \$1,076}
  \resumeItem{Calculated 823 lbf parachute shock and 1,230 lbf snatch loads using five published methods}
  \resumeItem{Developed CO$_2$ deployment hardware in a three-person team after two integrated-test failures, applying thick-wall stress and fastener calculations through CAD, machining, component testing, and a third full-system test}
  \resumeItem{Redesigned and machined cable cutters, reducing mass 49.5\% with a calculated 2.01 FOS}
  \resumeItem{Introduced Brummel eye splices to replace knots and quick-links; dynamically tested in-house splices on nonstandard line to a minimum 1.5 factor of safety (FOS)}
  \resumeItem{Wrote standard operating procedures (SOPs) for separation and recovery tests; led crew briefings, firing/abort decisions, and post-test analysis}
\resumeItemListEnd

\resumeSubheading
  {AIAA Design-Build-Fly Aircraft}
  {June 2025 - December 2025}
  {Wing Structural Lead, Boston University}
  {}
\resumeItemListStart
  \resumeItem{Led material, structural, and manufacturing trade studies for two Fusion 360 wing designs; produced an XPS prototype 50\% lighter than the prior year's wing and cut build time from three months to one}
  \resumeItem{Designed a split NACA 2412 wing and servo flap system in Fusion 360; mentored nine members in CAD and fabrication}
\resumeItemListEnd

\resumeSubheading
  {Hybrid Rocket Nozzle}
  {September 2024 - November 2024}
  {Lead Engineer, Boston University Rocket Propulsion Group}
  {}
\resumeItemListStart
  \resumeItem{Reduced nozzle mass 73\% using SolidWorks CAD/FEA; strengthened the contour after review identified handling and assembly fragility, calculating a 7.8 FOS under mechanical loads alone}
  \resumeItem{Led five members through three design reviews and a two-second hotfire}
\resumeItemListEnd

\resumeSubheading
  {Aerodynamic Downcomer Fairing}
  {August 2026 - Present}
  {Project Lead, Boston University Rocket Propulsion Group}
  {}
\resumeItemListStart
  \resumeItem{Leading set-based design of a fairing to protect Argo's downcomer from thermal and aerodynamic loads; comparing three concepts for supersonic drag, loading, thermal effects, and geometric constraints}
  \resumeItem{Developing Ansys CFD models at Mach 2.0--2.81; designing a curved carbon-fiber test article for combined thermal/static loading of bolted attachments (testing in development)}
\resumeItemListEnd

\resumeSubHeadingListEnd
\section{Technical Skills}
{\small
SolidWorks, Fusion 360, Onshape, FEA, Structural Load Analysis, Bolted Joints, GD\&T, Engineering Drawings, Tolerance Analysis, CNC \& Manual Machining, Composite Layup, Ansys CFD, Vacuum Systems, LabJack DAQ, Test Operations, Python, MATLAB
}
\end{document}
